% Clase 5 --- Convención de signo en superficies cerradas (§4.5).
% Para una superficie CERRADA la normal apunta siempre HACIA AFUERA: donde el
% campo entra el flujo es negativo y donde sale, positivo.
\begin{center}
\begin{tikzpicture}[scale=1]

  % campo uniforme de fondo
  \foreach \y in {-1.5,-0.75,0,0.75,1.5} {
    \draw[figvecaux] (-3.6,\y) -- (-2.6,\y);
  }
  \node[figlblaux, above=2pt] at (-3.1,1.55) {$\vec v$};

  % superficie cerrada
  \draw[figline, fill=figblue!6] (0,0) circle (1.6);

  % normales salientes
  \foreach \a in {0,45,...,315} {
    \draw[figvec, figgray] (\a:1.6) -- (\a:2.25);
  }
  \node[figlblaux, right=2pt] at (2.3,0) {$\hat n$ (siempre saliente)};

  % líneas de campo que atraviesan
  \foreach \y in {-1.0,-0.35,0.35,1.0} {
    \draw[figvecres] (-2.45,\y) -- (2.45,\y);
  }

  % rótulos de entrada y salida
  \node[figlbl, figblue, align=center] at (-2.35,-2.15)
    {entra: $\vec v \cdot \hat n < 0$\\ flujo \textbf{negativo}};
  \node[figlbl, figred, align=center] at (2.35,-2.15)
    {sale: $\vec v \cdot \hat n > 0$\\ flujo \textbf{positivo}};
  \node[figlbl] at (0,-3.0) {$\displaystyle \Phi = \oint_S \vec v \cdot d\vec A$};

\end{tikzpicture}\\[0.5em]
{\small\itshape Sin una convención, el signo de $\Phi$ sería arbitrario: una
superficie abierta admite dos normales opuestas y ambas son igual de válidas.
En una superficie \textbf{cerrada} el criterio ``$\hat n$ hacia afuera'' la fija
sin ambigüedad, y entonces el signo del flujo pasa a significar algo:
\textbf{sale más de lo que entra} o al revés.}
\end{center}
